\documentclass[11pt]{article}

\usepackage[margin=1in]{geometry}
\usepackage{amsmath,amssymb}
\usepackage{booktabs}
\usepackage{enumitem}
\usepackage{hyperref}
\usepackage{xcolor}

\hypersetup{
  colorlinks=true,
  linkcolor=blue!50!black,
  urlcolor=blue!50!black,
  citecolor=blue!50!black
}

\newcommand{\grip}{\texttt{grip}}
\newcommand{\Rfunc}[1]{\texttt{#1}}
\newcommand{\RR}{\mathbb{R}}
\newcommand{\norm}[1]{\left\lVert #1 \right\rVert}
\newcommand{\Eedge}{\mathcal{E}_{\mathrm{edge}}}
\newcommand{\Emetric}{\mathcal{E}_{\mathrm{metric}}}
\newcommand{\Erep}{\mathcal{E}_{\mathrm{rep}}}
\newcommand{\Eanchor}{\mathcal{E}_{\mathrm{anchor}}}

\title{Validation Plan for \texttt{grip.layout.misf.edge.kk()}}
\author{Design note for the \grip{} package}
\date{May 19, 2026}

\begin{document}

\maketitle

\begin{abstract}
This document translates the staged hybrid
\Rfunc{grip.layout.misf.edge.kk()} algorithm specification into concrete
validation gates. The goal is to make implementation auditable before writing
production code. The validation plan covers objective terms, finite-difference
gradients, active-set and constraint contracts, Armijo line search behavior,
scale policies, trace schema, equivalence limits, tuning reports, and release
readiness. The companion algorithm specification is
\texttt{dev/design/misf-edge-kk/misf\_edge\_kk\_algorithm\_2026-05-19.tex}.
\end{abstract}

\tableofcontents

\section{Purpose and Scope}

The proposed \Rfunc{grip.layout.misf.edge.kk()} algorithm is a staged hybrid:
\begin{enumerate}[leftmargin=2em]
  \item weighted-MISF construction;
  \item heuristic insertion;
  \item explicit Armijo objective refinement per active level;
  \item final pure or near-pure \Rfunc{grip.optimize.edge.kk.layout()} polish.
\end{enumerate}

This validation plan defines the gates that should be passed before the method
is treated as usable. It does not specify implementation details beyond what is
needed to make the algorithm testable.

\subsection{Primary Validation Goals}

\begin{itemize}[leftmargin=2em]
  \item Verify that each objective term is numerically correct.
  \item Verify that active sets and constraints match the MISF contract.
  \item Verify that Armijo refinement decreases the explicit level objective.
  \item Verify that degenerate settings reduce to known edge-KK behavior.
  \item Verify that trace output is rich enough to diagnose tuning and failures.
  \item Verify that schedule tuning can be compared against current baselines.
\end{itemize}

\subsection{Non-Goals}

The validation plan does not require the first implementation to be fastest,
fully adaptive, or better than all existing workflows on all graph families. The
first target is correctness, auditability, and diagnostic completeness.

\section{Gate Naming Convention}

Each gate has an identifier:
\[
  \texttt{MEK-<area>-<number>}.
\]
The areas are:
\begin{description}[leftmargin=2em]
  \item[OBJ] objective values and gradients;
  \item[MISF] active-set and hierarchy contract;
  \item[CON] constraint construction;
  \item[ARM] Armijo optimizer behavior;
  \item[SCL] scale policy;
  \item[TRC] trace schema;
  \item[EQV] equivalence and degenerate cases;
  \item[TUN] schedule tuning and benchmark reports;
  \item[REL] release readiness.
\end{description}

\section{Reference Objective}

At active level \(q\), the implementation should refine only
\[
  Z_{A_q} = \{z_i : i\in A_q\}
\]
under
\[
  \mathcal{E}_q(Z_{A_q})
  =
  \mathcal{E}_{\mathrm{edge},q}(Z_{A_q})
  +
  \rho_q \mathcal{E}_{\mathrm{metric},q}(Z_{A_q})
  +
  \lambda_q \mathcal{E}_{\mathrm{rep},q}(Z_{A_q})
  +
  \mathcal{E}_{\mathrm{anchor},q}(Z_{A_q}).
\]

The true edge term uses only original graph edges:
\[
  E_{\mathrm{active},q}
  =
  \{(i,j)\in E : i,j\in A_q\}.
\]

The auxiliary metric term uses temporary weighted graph-geodesic or
metric-neighbor constraints:
\[
  M_q =
  \{(i,j,\delta_{ij}) : i,j\in A_q,\ i<j,\ \delta_{ij}=d_G(i,j)\},
\]
deduplicated and disjoint from \(E_{\mathrm{active},q}\).

The repulsion term uses active pairs \(P_q\subseteq A_q\times A_q\), exact below
a threshold and deterministic sampled above it. The anchor term uses coordinates
immediately after insertion and before level refinement.

\section{Objective and Gradient Gates}

\subsection{MEK-OBJ-001: Edge-Stress Value}

For a fixed small graph, fixed coordinates, fixed edge scale, and known
stiffnesses, the implementation must compute
\[
  \mathcal{E}_{\mathrm{edge},q}(Z)
  =
  \frac{1}{2}
  \sum_{(i,j)\in E_{\mathrm{active},q}}
  k_{ij}^{(q)}
  \left(r_{ij}(Z)-s_q\ell_{ij}\right)^2
\]
to within numerical tolerance of a simple R reference implementation.

\textbf{Suggested tolerance:}
\[
  |\widehat{\Eedge}-\Eedge^{\mathrm{ref}}|
  \le 10^{-10}\max(1,|\Eedge^{\mathrm{ref}}|).
\]

\subsection{MEK-OBJ-002: Edge-Stress Gradient}

The edge-stress gradient must match centered finite differences on random
small fixtures in dimensions 2, 3, and 5.

For coordinate \(x_k\), use
\[
  \frac{\mathcal{E}(x_k+h)-\mathcal{E}(x_k-h)}{2h}
\]
with \(h\in[10^{-6},10^{-5}]\).

\textbf{Suggested tolerance:}
\[
  \frac{\norm{G-G_{\mathrm{fd}}}_F}{\max(1,\norm{G_{\mathrm{fd}}}_F)}
  \le 10^{-5}.
\]

\subsection{MEK-OBJ-003: Auxiliary Metric-Stress Value and Gradient}

Repeat MEK-OBJ-001 and MEK-OBJ-002 for
\(\mathcal{E}_{\mathrm{metric},q}\), using metric-neighbor
constraints \(M_q\). This test must explicitly include cases where an original
edge appears in the metric-neighbor cache; it must be removed from \(M_q\) and
kept only in \(E_{\mathrm{active},q}\).

\subsection{MEK-OBJ-004: Log-Repulsion Value and Gradient}

For fixed \(P_q\) and fixed pair weights, verify
\[
  \mathcal{E}_{\mathrm{rep},q}(Z)
  =
  -
  \sum_{(i,j)\in P_q}
  \omega_{ij}^{(q)}
  \log\left(
    \sqrt{\norm{z_i-z_j}^2+\epsilon_{\mathrm{rep}}^2}
  \right)
\]
against an R reference value and finite-difference gradient.

The test panel must include:
\begin{itemize}[leftmargin=2em]
  \item near-colliding vertices;
  \item exactly repeated coordinates with positive \(\epsilon_{\mathrm{rep}}\);
  \item dimensions 2, 3, and 5;
  \item exact and sampled \(P_q\).
\end{itemize}

\subsection{MEK-OBJ-005: Anchor Value and Gradient}

For fixed anchor targets \(\tilde{Z}^{(q)}\) and weights \(a_i^{(q)}\), verify
\[
  \mathcal{E}_{\mathrm{anchor},q}(Z)
  =
  \frac{1}{2}
  \sum_{i\in A_q}
  a_i^{(q)}
  \norm{z_i-\tilde{z}_i^{(q)}}^2
\]
and its gradient. Include tests with zero anchor weights and different weights
for newly inserted vertices.

\subsection{MEK-OBJ-006: Total Objective Composition}

The total objective and gradient must equal the weighted sum of independently
computed term objectives and gradients:
\[
  \nabla \mathcal{E}_q
  =
  \nabla\mathcal{E}_{\mathrm{edge},q}
  +
  \rho_q\nabla\mathcal{E}_{\mathrm{metric},q}
  +
  \lambda_q\nabla\mathcal{E}_{\mathrm{rep},q}
  +
  \nabla\mathcal{E}_{\mathrm{anchor},q}.
\]

The trace must report the same term energies that sum to the total objective.

\section{MISF Active-Set Gates}

\subsection{MEK-MISF-001: Prefix Active Sets}

For every level \(q\), the active set must be exactly
\[
  A_q = \{\pi_1,\ldots,\pi_{m_q}\},
\]
where \(\pi=\texttt{misf.order}\) and \(m_q=\texttt{misf.level\_size[q]}\).

The test should construct a small known MISF object and verify exact equality
of active vertex IDs and order.

\subsection{MEK-MISF-002: Newly Inserted Vertices}

For every transition \(q+1\to q\), verify
\[
  N_q = A_q\setminus A_{q+1}.
\]
The inserted count in the trace must equal \(|N_q|\).

\subsection{MEK-MISF-003: Inactive Vertices Excluded}

Vertices outside \(A_q\) must not appear in:
\begin{itemize}[leftmargin=2em]
  \item \(E_{\mathrm{active},q}\);
  \item \(M_q\);
  \item \(P_q\);
  \item anchor records;
  \item active-level gradient rows.
\end{itemize}

Inactive coordinates may exist in storage, but they must not affect the
active-level objective.

\subsection{MEK-MISF-004: Previously Active Vertices Can Move}

During Armijo refinement at level \(q\), both \(A_{q+1}\) and \(N_q\) are
optimization variables. A test should construct a case where an already-active
vertex has nonzero gradient and confirm it is allowed to move after an accepted
step.

\section{Constraint Construction Gates}

\subsection{MEK-CON-001: Original Active Edge Set}

Verify that
\[
  E_{\mathrm{active},q}
  =
  \{(i,j)\in E : i,j\in A_q\}
\]
contains all and only original graph edges with both endpoints active.

Tests should include:
\begin{itemize}[leftmargin=2em]
  \item an edge with one endpoint inactive;
  \item an edge with both endpoints active;
  \item multiple active levels;
  \item weighted edges with nonunit target lengths.
\end{itemize}

\subsection{MEK-CON-002: Metric Constraints Are Separate}

The temporary metric set \(M_q\) must be deduplicated as unordered pairs and
disjoint from \(E_{\mathrm{active},q}\). If a metric-neighbor pair is also an
original edge, it belongs only to \(E_{\mathrm{active},q}\).

\subsection{MEK-CON-003: Metric Targets Are Graph-Geodesic Distances}

Every \((i,j,\delta_{ij})\in M_q\) must use the weighted graph-geodesic distance
\(d_G(i,j)\) under the normalized edge lengths used by the solver.

For small graphs, compare against an independent all-pairs Dijkstra reference.

\subsection{MEK-CON-004: Repulsion Pair Mode}

Below the exact threshold, \(P_q\) must contain all unordered active pairs.
Above the threshold, \(P_q\) must be a deterministic sample keyed by graph,
seed, level, and active order.

\subsection{MEK-CON-005: Repulsion Weights Normalize Pair Mass}

For each level,
\[
  \sum_{(i,j)\in P_q}\omega_{ij}^{(q)} = |A_q|
\]
or another explicitly documented normalization. The chosen normalization must
be reported in the trace.

\subsection{MEK-CON-006: Anchor Targets Are Post-Insertion Coordinates}

Anchor targets must be captured after heuristic insertion and before
active-level refinement. A test should perturb coordinates during refinement
and confirm anchors still point to the stored post-insertion coordinates, not
to moving coordinates.

\section{Armijo Optimizer Gates}

\subsection{MEK-ARM-001: Fixed Objective During Line Search}

During a single Armijo line search, \(E_{\mathrm{active},q}\), \(M_q\), \(P_q\),
pair weights, anchor targets, \(\rho_q\), \(\lambda_q\), anchor weights, and
scale must remain fixed. This is especially important for sampled \(P_q\).

\subsection{MEK-ARM-002: Objective Decrease}

Accepted steps must satisfy
\[
  \mathcal{E}_q(Z')
  \le
  \mathcal{E}_q(Z)
  -
  c\eta\norm{G_q}_F^2.
\]

Tests should run small levels and assert monotone decrease across accepted
steps.

\subsection{MEK-ARM-003: Rejected Step Shrinkage}

When a trial step fails the Armijo condition, the next trial step must shrink
by the configured factor. The trace must record the accepted step and number of
shrinks.

\subsection{MEK-ARM-004: Active Recenter}

After accepted steps, active coordinates should be recentered. The test should
verify
\[
  \sum_{i\in A_q} z_i = 0
\]
to numerical tolerance after accepted steps, while inactive vertices remain
irrelevant to the active-level center.

\subsection{MEK-ARM-005: Stopping Reasons}

Each level trace must report one of:
\begin{itemize}[leftmargin=2em]
  \item iteration budget reached;
  \item gradient tolerance reached;
  \item minimum step reached;
  \item relative improvement patience reached;
  \item no accepted Armijo step.
\end{itemize}

\subsection{MEK-ARM-006: Finite Output Under Stress Cases}

The optimizer must return finite coordinates and finite trace values for:
\begin{itemize}[leftmargin=2em]
  \item repeated initial coordinates;
  \item zero metric coefficient \(\rho_q=0\);
  \item zero repulsion coefficient \(\lambda_q=0\);
  \item zero anchor weights;
  \item sparse active edge sets.
\end{itemize}

\section{Scale Policy Gates}

\subsection{MEK-SCL-001: No Profiled Scale With Active Repulsion}

During MISF refinement, profiled scale is forbidden whenever
\[
  \lambda_q > 0 \quad \text{or} \quad \rho_q > 0.
\]
Attempting this combination should error or downgrade only through an explicit
documented choice. Silent profiling is not allowed.

\subsection{MEK-SCL-002: Identity Scale}

With identity scale, the target length for a true edge is exactly \(\ell_{ij}\).
This mode should be the simplest deterministic baseline.

\subsection{MEK-SCL-003: Global Fixed Scale}

With global fixed scale, the same scale must be used at every MISF level. The
trace must report the same scale value at all levels.

\subsection{MEK-SCL-004: Level Fixed Scale}

With level fixed scale, the scale may differ by level but must remain fixed
within all Armijo iterations of that level. This mode is diagnostic and should
be trace-visible.

\subsection{MEK-SCL-005: Final Profiled Polish}

Profiled scale is allowed for the final edge-KK polish when repulsion and
auxiliary metric terms are off. The final polish scale and final edge RMSE must
be reported.

\section{Trace Schema Gates}

\subsection{MEK-TRC-001: Required Per-Level Fields}

Each level summary must include:
\begin{itemize}[leftmargin=2em]
  \item level index;
  \item active count \(|A_q|\);
  \item inserted count \(|N_q|\);
  \item counts of \(E_{\mathrm{active},q}\), \(M_q\), and \(P_q\);
  \item \(\rho_q\), \(\lambda_q\), and anchor-weight summaries;
  \item scale mode and scale value;
  \item stopping reason;
  \item elapsed time.
\end{itemize}

\subsection{MEK-TRC-002: Required Per-Iteration Fields}

Each Armijo iteration row must include:
\begin{itemize}[leftmargin=2em]
  \item total objective;
  \item edge, metric, repulsion, and anchor energies;
  \item gradient norm;
  \item accepted step size;
  \item number of shrink attempts;
  \item true-edge relative RMSE;
  \item auxiliary metric relative RMSE when \(|M_q|>0\);
  \item elapsed time.
\end{itemize}

\subsection{MEK-TRC-003: Term Energies Sum to Total}

For every trace row,
\[
  \mathcal{E}_{\mathrm{total}}
  =
  \mathcal{E}_{\mathrm{edge}}
  +
  \rho_q\mathcal{E}_{\mathrm{metric}}
  +
  \lambda_q\mathcal{E}_{\mathrm{rep}}
  +
  \mathcal{E}_{\mathrm{anchor}}
\]
to numerical tolerance.

\subsection{MEK-TRC-004: Adaptive Decisions Are Trace-Visible}

If adaptive guards are implemented later, every adaptive change must record:
\begin{itemize}[leftmargin=2em]
  \item trigger diagnostic;
  \item previous value;
  \item new value;
  \item level and iteration.
\end{itemize}

\section{Equivalence and Degenerate Gates}

\subsection{MEK-EQV-001: One-Level Edge-KK Equivalence}

When the solver is configured with:
\begin{itemize}[leftmargin=2em]
  \item one active level;
  \item \(\rho_q=0\);
  \item \(\lambda_q=0\);
  \item anchors off;
  \item the same scale mode and stiffnesses as
  \Rfunc{grip.optimize.edge.kk.layout()};
\end{itemize}
the result should match \Rfunc{grip.optimize.edge.kk.layout()} within numerical
tolerance, modulo equivalent trace formatting.

This is the most important correctness bridge to the existing edge-KK code.

\subsection{MEK-EQV-002: Zero Auxiliary Metric}

With \(\rho_q=0\), the metric constraint set may still be constructed for trace
diagnostics, but it must not affect objective, gradient, Armijo decisions, or
coordinates.

\subsection{MEK-EQV-003: Zero Repulsion}

With \(\lambda_q=0\), the repulsion pair set may still be constructed for trace
diagnostics, but it must not affect objective, gradient, Armijo decisions, or
coordinates.

\subsection{MEK-EQV-004: Zero Anchors}

With all anchor weights set to zero, anchor targets must not affect objective,
gradient, Armijo decisions, or coordinates.

\subsection{MEK-EQV-005: Dimension Support}

The solver must produce finite deterministic outputs for dimensions 2, 3, 4,
and at least one larger dimension such as 8 on small graphs.

\section{Tuning and Benchmark Gates}

\subsection{MEK-TUN-001: Fixture Matrix}

The tuning report must include at least:
\begin{itemize}[leftmargin=2em]
  \item paths;
  \item cycles;
  \item flat 2D grids;
  \item nonuniform 2D grids;
  \item flat 3D lattices;
  \item nonuniform 3D grids;
  \item 2D and 3D quadform-derived fixtures where available;
  \item trees;
  \item irregular manifold-like examples;
  \item sparse bottleneck or density-varying graphs.
\end{itemize}

\subsection{MEK-TUN-002: Baseline Matrix}

Every fixture should be compared against:
\begin{enumerate}[leftmargin=2em]
  \item \Rfunc{grip.layout.weighted()} followed by
  \Rfunc{grip.optimize.edge.kk.layout()};
  \item \Rfunc{grip.metric.mds.layout()} followed by
  \Rfunc{grip.optimize.edge.kk.layout()};
  \item \Rfunc{grip.layout.misf.edge.kk()} with candidate schedules;
  \item a zero-regularization MISF diagnostic where meaningful.
\end{enumerate}

\subsection{MEK-TUN-003: Schedule Grid}

The first grid should cover a small family:
\[
  \rho_q = \rho_{\max}\theta_q^{\gamma_\rho},
  \qquad
  \lambda_q = \lambda_{\min}+\lambda_{\max}\theta_q^{\gamma_\lambda},
\]
with a limited set of \(\rho_{\max}\), \(\lambda_{\max}\), and exponent values.
The exact grid can be revised after objective magnitudes are observed, but the
report must record all tested values.

\subsection{MEK-TUN-004: Metrics}

The tuning report must include:
\begin{itemize}[leftmargin=2em]
  \item final true-edge relative RMSE;
  \item true-edge signed bias;
  \item graph-geodesic distortion;
  \item auxiliary metric RMSE where active;
  \item local angle or quadform preservation where available;
  \item Procrustes RMSE to known coordinates where available;
  \item fold-over or cell-orientation failures for grid-like fixtures;
  \item seed stability;
  \item Armijo rejected-step counts;
  \item runtime.
\end{itemize}

\subsection{MEK-TUN-005: Pareto Selection}

The default schedule should be selected as a Pareto compromise. A candidate
should not be chosen solely because it minimizes edge RMSE if it causes
fold-over, instability, or excessive runtime.

\subsection{MEK-TUN-006: Adaptive Guards}

Adaptive guards are optional for the first implementation. If included, each
guard must have a direct diagnostic trigger, for example:
\begin{itemize}[leftmargin=2em]
  \item active-edge graph too sparse: increase \(\rho_q\);
  \item active pair distances collapsing: increase \(\lambda_q\);
  \item repulsion dominates late: decrease \(\lambda_q\);
  \item metric stress dominates late: decrease \(\rho_q\);
  \item repeated Armijo failure: reduce regularization before reducing edge
  objective.
\end{itemize}

Each adjustment must be trace-visible.

\section{Report Artifacts}

\subsection{Validation Report}

A validation run should produce a report under a path such as:
\[
  \texttt{output/misf-edge-kk-validation/}
\]
with:
\begin{itemize}[leftmargin=2em]
  \item objective-gradient finite-difference tables;
  \item active-set and constraint audit tables;
  \item Armijo decrease summaries;
  \item equivalence test summaries;
  \item trace schema checks;
  \item tuning benchmark summaries.
\end{itemize}

\subsection{Machine-Readable Outputs}

The report should write CSV or RDS files for:
\begin{itemize}[leftmargin=2em]
  \item per-test gate status;
  \item per-level trace summaries;
  \item per-iteration trace summaries;
  \item per-fixture benchmark metrics;
  \item schedule grid results.
\end{itemize}

\section{Release Readiness Gates}

\subsection{MEK-REL-001: Internal Prototype Readiness}

The method is ready for an internal R prototype when:
\begin{itemize}[leftmargin=2em]
  \item the objective and gradient tests pass for R reference functions;
  \item active-set and constraint construction tests pass;
  \item one-level edge-KK equivalence is specified and testable;
  \item trace schema is implemented at least in R.
\end{itemize}

\subsection{MEK-REL-002: Compiled Backend Readiness}

The method is ready for compiled implementation when:
\begin{itemize}[leftmargin=2em]
  \item the R prototype passes all objective and Armijo tests;
  \item compiled objective values and gradients match the R reference;
  \item compiled traces match R prototype traces on tiny fixtures.
\end{itemize}

\subsection{MEK-REL-003: Experimental Export Readiness}

The method is ready for an experimental exported wrapper when:
\begin{itemize}[leftmargin=2em]
  \item all objective, MISF, constraint, Armijo, scale, trace, and equivalence
  gates pass;
  \item dimension 2, 3, 4, and high-dimensional smoke tests pass;
  \item tuning reports identify at least one reasonable global default;
  \item current pipelines remain available and documented as baselines;
  \item the function is clearly marked experimental.
\end{itemize}

\subsection{MEK-REL-004: Production Readiness}

The method should not be considered production-ready until:
\begin{itemize}[leftmargin=2em]
  \item the tuning panel shows robust behavior across graph families;
  \item final edge-KK polish is consistently competitive with existing
  workflows;
  \item failures are trace-diagnosable;
  \item runtime is acceptable for target graph sizes;
  \item package checks remain stable.
\end{itemize}

\section{Minimal First Validation Milestone}

The first practical milestone should validate a one-component R prototype on
small graphs. The milestone passes when:
\begin{enumerate}[leftmargin=2em]
  \item edge, metric, repulsion, and anchor gradients pass finite differences;
  \item active sets are exact MISF prefixes;
  \item \(M_q\) is disjoint from \(E_{\mathrm{active},q}\);
  \item \(P_q\) is deterministic and normalized;
  \item Armijo accepted steps decrease the total objective;
  \item the trace contains all required fields;
  \item the one-level, zero-regularization case matches
  \Rfunc{grip.optimize.edge.kk.layout()}.
\end{enumerate}

\section{Conclusion}

This validation plan makes \Rfunc{grip.layout.misf.edge.kk()} testable before
implementation. The key release gate is not visual plausibility, but objective
correctness under a transparent staged hybrid contract: true edge-KK stress,
separate temporary metric-neighbor stress, normalized log repulsion, quadratic
anchors, Armijo descent, and final pure or near-pure edge-KK polish.

\end{document}
